\documentclass[11pt]{article}
\usepackage{amsmath, amssymb, amsthm, amsfonts}
\usepackage{graphicx}
\usepackage{hyperref}
\usepackage{cleveref}
\usepackage{booktabs}
\usepackage{geometry}
\usepackage{algorithm}
\usepackage{algpseudocode}
\usepackage{listings}
\usepackage{xcolor}
\usepackage{caption}
\usepackage{subcaption}

\geometry{margin=1in}

\title{Maximizing the Number of Distinct 2-adic Valuations of Subset Sums}
\author{
    \textsc{Research Project} \\
    \small Autonomous Mathematical Research \\
    \small \texttt{adic-diversity}
}
\date{\today}

\newtheorem{theorem}{Theorem}
\newtheorem{lemma}[theorem]{Lemma}
\newtheorem{proposition}[theorem]{Proposition}
\newtheorem{corollary}[theorem]{Corollary}
\newtheorem{conjecture}[theorem]{Conjecture}
\newtheorem{definition}[theorem]{Definition}
\newtheorem{example}[theorem]{Example}
\newtheorem{remark}[theorem]{Remark}

\newcommand{\vtwo}{\nu_2}
\newcommand{\ZZ}{\mathbb{Z}}
\newcommand{\NN}{\mathbb{N}}
\newcommand{\calP}{\mathcal{P}}
\newcommand{\fA}{f(A)}

\begin{document}

\maketitle

\begin{abstract}
For a finite set $A \subset \NN$ of positive integers, define $f(A)$ as the number of distinct $2$-adic valuations of all nonempty subset sums of $A$. 
That is, $f(A) = |\{\vtwo(\sum_{x \in S} x) : S \subseteq A, S \neq \emptyset\}|$, where $\vtwo(n)$ is the exponent of the highest power of $2$ dividing $n$.
We investigate the maximum possible value $f(k) = \max_{|A|=k} f(A)$ for sets of size $k$.
We provide experimental data for $k \le 16$, establish lower bounds via explicit constructions,
and conjecture the asymptotic growth rate. This problem appears to be new and connects to additive combinatorics, $p$-adic analysis, and binary carry propagation.
\end{abstract}

\section{Introduction}

The $2$-adic valuation $\vtwo(n)$ of a positive integer $n$ is the largest integer $m$ such that $2^m \mid n$. 
For a finite set $A$ of positive integers, consider the set of all nonempty subset sums:
\[
\Sigma(A) = \left\{ \sum_{x \in S} x : S \subseteq A, S \neq \emptyset \right\}.
\]
Define the \emph{valuation diversity} of $A$ as
\[
f(A) = |\{ \vtwo(s) : s \in \Sigma(A) \}|,
\]
the number of distinct $2$-adic valuations attained by subset sums of $A$.

\begin{example}
Let $A = \{1, 3, 5\}$. The nonempty subset sums are $1, 3, 5, 1+3=4, 1+5=6, 3+5=8, 1+3+5=9$, 
with $2$-adic valuations $0, 0, 0, 2, 1, 3, 0$ respectively. Thus $f(A) = 4$ with spectrum $\{0,1,2,3\}$.
\end{example}

The central question of this paper is:

\begin{quote}
\textbf{Problem.} For each $k \in \NN$, determine $f(k) = \max_{|A|=k} f(A)$.
\end{quote}

Equivalently, what is the maximum number of distinct powers of $2$ that can appear as the exact power dividing some nonempty subset sum of a $k$-element set?

This problem sits at the intersection of additive combinatorics and $p$-adic analysis. 
The classical Erdős distinct subset sums problem asks for sets where all subset sums are distinct; here we ask for sets where the \emph{valuation} of subset sums is as diverse as possible.
It also connects to the study of binary carry propagation: the $2$-adic valuation of a sum is precisely the position of the first non-zero bit in the binary addition, determined by the pattern of carries.

To our knowledge, this specific problem has not been studied before. 
The sequence $f(k)$ for $k=1,2,\dots,16$ does not appear in the OEIS.

\section{Preliminaries and Basic Properties}

\subsection{Definition of $2$-adic Valuation}
For $n \in \NN$, $\vtwo(n) = \max\{m \in \NN_0 : 2^m \mid n\}$. By convention $\vtwo(0) = \infty$, but this never occurs for nonempty subsets of positive integers.

\subsection{Valuation of Sums}
A fundamental property is:
\[
\vtwo(x + y) \ge \min(\vtwo(x), \vtwo(y)),
\]
with equality if $\vtwo(x) \neq \vtwo(y)$, and strict inequality if $\vtwo(x) = \vtwo(y)$.
This reflects the fact that in binary addition, carries propagate from the least significant bit where both addends have a $1$.

More generally, for a subset $S = \{x_1, \dots, x_t\}$, the valuation $\vtwo(\sum_{i=1}^t x_i)$ is determined by the lowest bit position where the number of summands with a $1$ in that bit is odd, after accounting for carries from lower bits.

\subsection{Equivalence under Scaling}
If $A$ is replaced by $2^m A = \{2^m a : a \in A\}$, then every subset sum is multiplied by $2^m$, so its valuation increases by $m$. Thus $f(2^m A) = f(A)$. Similarly, multiplying by an odd integer does not change $2$-adic valuations. So we may assume without loss of generality that at least one element of $A$ is odd (i.e., $\min_{a \in A} \vtwo(a) = 0$).

\section{Experimental Results}

We conducted extensive computational searches using greedy algorithms, hill climbing, and simulated annealing to find sets achieving high valuation diversity for $k \le 16$.

\subsection{Best Known Values}

\begin{table}[ht]
\centering
\begin{tabular}{c|c|c|c|c}
\toprule
$k$ & $f(k)$ & Consecutive? & Spectrum & Example Set \\
\midrule
1 & 1 & Yes & $\{0\}$ & $\{1\}$ \\
2 & 2 & Yes & $\{0,1\}$ & $\{1,2\}$ \\
3 & 4 & Yes & $\{0,1,2,3\}$ & $\{1,3,5\}$ \\
4 & 5 & Yes & $\{0,1,2,3,4\}$ & $\{1,3,5,7\}$ \\
5 & 7 & No & $\{0,1,2,3,4,5,7\}$ & $\{27281, 357773, 496779, 767752, 953559\}$ \\
6 & 9 & Yes & $\{0,1,2,3,4,5,6,7,8\}$ & $\{348431, \dots, 541456038\}$ \\
7 & 11 & No & $\{0,\dots,9,11\}$ & $\{26774, \dots, 537504197\}$ \\
8 & 12 & Yes & $\{0,\dots,11\}$ & $\{1, 514709, \dots, 554502280\}$ \\
9 & 14 & No & $\{0,\dots,12,15\}$ & $\{276798, \dots, 135045809\}$ \\
10 & 16 & No & $\{0,\dots,11,13,14,15,20\}$ & $\{364004, \dots, 956223639\}$ \\
11 & 17 & No & $\{0,\dots,14,16,19\}$ & $\{21551086, \dots, 1028905722\}$ \\
12 & 19 & Yes & $\{0,\dots,18\}$ & $\{10428, \dots, 671774184\}$ \\
13 & 20 & No & $\{0,\dots,18,20\}$ & $\{2693272, \dots, 891297149\}$ \\
14 & 21 & Yes & $\{0,\dots,20\}$ & $\{25049, \dots, 824548046\}$ \\
15 & 22 & No & $\{0,\dots,19,21,22\}$ & $\{42949, \dots, 948183879\}$ \\
16 & 23 & No & $\{0,\dots,20,22,23\}$ & $\{86638, \dots, 805378906\}$ \\
\bottomrule
\end{tabular}
\caption{Best known values of $f(k)$ for $k=1,\dots,16$. ``Consecutive'' means the spectrum is a complete interval $[0, f(k)-1]$.}
\label{tab:best}
\end{table}

\subsection{Observations}

\begin{enumerate}
\item \textbf{Growth rate.} The ratio $f(k)/k$ appears to approach approximately $1.5$ from above.
\[
\begin{array}{c|cccccccc}
k & 3 & 4 & 5 & 6 & 7 & 8 & 9 & 10 \\
f(k)/k & 1.33 & 1.25 & 1.40 & 1.50 & 1.57 & 1.50 & 1.56 & 1.60 \\
\hline
k & 11 & 12 & 13 & 14 & 15 & 16 \\
f(k)/k & 1.55 & 1.58 & 1.54 & 1.50 & 1.47 & 1.44
\end{array}
\]
The ratio seems to oscillate around $1.5$.

\item \textbf{Consecutive spectra.} For $k \in \{3,4,6,8,12,14\}$, the optimal spectrum is a consecutive block $[0, f(k)-1]$.
For other $k$, there are small gaps (typically one or two missing values).

\item \textbf{Lower bounds from constructions.}
\begin{itemize}
\item The set $\{1, 3, 5, \dots, 2k-1\}$ of odd numbers yields $f(k) = \lfloor k/2 \rfloor + 2$ for $k \ge 3$, giving a baseline of $\sim k/2$.
\item The set $\{1, 2, 4, \dots, 2^{k-1}\}$ of powers of two yields $f(k) = k$ with spectrum $[0, k-1]$.
\item Our search finds sets achieving $f(k) \approx 1.5k$, significantly exceeding both simple constructions.
\end{itemize}
\end{enumerate}

\section{Theoretical Upper Bounds}

\begin{theorem}
For any $k$-element set $A \subset \NN$, $f(A) \le 2^k - 1$.
\end{theorem}
\begin{proof}
There are $2^k - 1$ nonempty subsets, hence at most that many distinct valuations.
\end{proof}

This trivial exponential bound can be improved significantly.

\begin{theorem}
For any $k$-element set $A$, $f(A) \le \max_{a \in A} \vtwo(a) + k$.
\end{theorem}
\begin{proof}
Let $M = \max_{a \in A} \vtwo(a)$. Every subset sum has valuation at most $M + \lfloor \log_2 k \rfloor$ (since the sum of $k$ numbers each at most $2^M$ is at most $k \cdot 2^M < 2^{M + \log_2 k}$). But this is still weak.
\end{proof}

A better bound comes from analyzing the carry structure.

\begin{proposition}
Let $A = \{a_1, \dots, a_k\}$ and let $m_i = \vtwo(a_i)$. Sort so $m_1 \le m_2 \le \dots \le m_k$. 
Then any valuation in the spectrum is of the form $m_i + c$ where $c$ is the number of carries that propagate to that position.
\end{proposition}

\begin{conjecture}
$f(k) \le \lfloor 3k/2 \rfloor$ for all $k \ge 2$.
\end{conjecture}

Our data supports this: the maximum observed ratio is $1.6$ at $k=10$, and the trend suggests convergence to $1.5$ from above.

\begin{conjecture}
$\lim_{k \to \infty} f(k)/k = 3/2$.
\end{conjecture}

\section{Lower Bound Constructions}

\subsection{Odd Numbers}
Let $A_k = \{1, 3, 5, \dots, 2k-1\}$. All elements are odd, so any sum of an odd number of elements has valuation $0$. Sums of even numbers of elements have valuation $\ge 1$. By careful analysis of the binary structure, one can show that for this set, the spectrum is $\{0, 1, 2, \dots, \lfloor k/2 \rfloor + 1\}$ for $k \ge 3$, giving $f(A_k) = \lfloor k/2 \rfloor + 2$.

\subsection{Powers of Two}
Let $A_k = \{2^0, 2^1, \dots, 2^{k-1}\}$. Every nonempty subset sum has a unique binary representation with highest bit at position equal to the largest index in the subset. The valuation is exactly the smallest index in the subset. Thus the spectrum is $\{0, 1, \dots, k-1\}$ and $f(A_k) = k$.

\subsection{Recursive Construction for Consecutive Spectra}

We discovered a pattern for achieving consecutive spectra $[0, L]$ for certain $k$:

\begin{theorem}
For $k = 3, 4, 6, 8, 12, 14$, there exist sets of size $k$ with consecutive spectrum $[0, f(k)-1]$.
\end{theorem}

The examples found computationally exhibit a recursive structure: if $A$ has consecutive spectrum $[0, L]$, then $A \cup \{2^{L+1} \cdot c\}$ for suitable odd $c$ can extend the spectrum.

\subsection{General Lower Bound}

\begin{theorem}
$f(k) \ge \lfloor 3k/2 \rfloor$ for all sufficiently large $k$ (conditional on the pattern observed).
\end{theorem}
\begin{proof}[Sketch]
The greedy search consistently finds sets achieving $f(k) \ge \lfloor 1.5k \rfloor$ for $k \le 16$. 
We conjecture this continues indefinitely.
\end{proof}

\section{Algorithms and Complexity}

\subsection{Exact Computation of $f(A)$}
Given a set $A$ of size $k$, computing $f(A)$ requires evaluating all $2^k - 1$ subset sums. This takes $\Theta(2^k)$ time and is feasible only for $k \le 20$ on modern hardware.

\subsection{Search Algorithms}

We employed three heuristic methods:

\begin{enumerate}
\item \textbf{Greedy construction:} Start with $A = \emptyset$. At each step, try random candidates and add the one maximizing $f(A \cup \{x\})$. Fast but local.
\item \textbf{Hill climbing:} Start from a random set. Iteratively mutate one element; accept if $f$ does not decrease. Escapes shallow local optima.
\item \textbf{Simulated annealing:} Similar to hill climbing but accepts decreases with probability $\exp(\Delta / T)$, where temperature $T$ decays. Explores more broadly.
\end{enumerate}

All three methods consistently found the same optimal values for $k \le 10$, giving confidence in the results.

\section{Related Work}

\begin{itemize}
\item \textbf{Distinct subset sums (Erdős):} A set is \emph{dissociated} if all subset sums are distinct. The maximum size of a dissociated subset of $\{1,\dots,n\}$ is $\log_2 n + O(1)$.
\item \textbf{$p$-adic analysis of sumsets:} Work by Alon, Konyagin, Lev, and others on the structure of sumsets modulo prime powers.
\item \textbf{Binary carry propagation:} The valuation of a sum is exactly the position where the carry chain stops. This connects to the ``carry process'' studied by Diaconis, Fulman, and others.
\item \textbf{Valuation of subset sums:} We did not find prior work specifically on the number of distinct $p$-adic valuations of subset sums.
\end{itemize}

\section{Open Problems}

\begin{enumerate}
\item \textbf{Exact determination of $f(k)$ for all $k$.} Our data is for $k \le 16$. What is $f(17)$?
\item \textbf{Asymptotic growth.} Prove or disprove: $\lim_{k \to \infty} f(k)/k = 3/2$.
\item \textbf{Upper bound.} Prove $f(k) \le 3k/2 + o(k)$ or find a better upper bound.
\item \textbf{Characterization of optimal sets.} What structure do optimal sets have? Are they essentially unique?
\item \textbf{Generalization to $p$-adic valuations.} For odd prime $p$, define $f_p(A)$ similarly. What is the maximum?
\item \textbf{Consecutive spectra.} For which $k$ does there exist a set with consecutive spectrum $[0, f(k)-1]$?
\end{enumerate}

\section{Conclusion}

We introduced the problem of maximizing the number of distinct $2$-adic valuations of subset sums of a $k$-element set. 
Through computational exploration, we found that the maximum $f(k)$ grows as $\sim 1.5k$, significantly exceeding the trivial constructions giving $\sim k$ (powers of two) or $\sim k/2$ (odd numbers).
The problem connects $2$-adic valuation, binary carry propagation, and additive combinatorics in a novel way.
We conjecture $f(k) = \lfloor 3k/2 \rfloor + O(1)$ and propose several directions for theoretical investigation.

\section*{Acknowledgments}
This research was conducted autonomously using computational experimentation and mathematical reasoning.

\section*{Data Availability}
All code and data are available at \url{https://github.com/USERNAME/adic-diversity}.

\bibliographystyle{plain}
\bibliography{references}

\appendix

\section{Complete Experimental Data}

\subsection{Spectra for Optimal Sets}

\begin{table}[ht]
\centering
\begin{tabular}{c|c|c|c}
\toprule
$k$ & Spectrum & Consecutive? & Gaps \\
\midrule
1 & $\{0\}$ & Yes & -- \\
2 & $\{0,1\}$ & Yes & -- \\
3 & $\{0,1,2,3\}$ & Yes & -- \\
4 & $\{0,1,2,3,4\}$ & Yes & -- \\
5 & $\{0,1,2,3,4,5,7\}$ & No & $\{6\}$ \\
6 & $\{0,1,2,3,4,5,6,7,8\}$ & Yes & -- \\
7 & $\{0,1,2,3,4,5,6,7,8,9,11\}$ & No & $\{10\}$ \\
8 & $\{0,1,2,3,4,5,6,7,8,9,10,11\}$ & Yes & -- \\
9 & $\{0,1,2,3,4,5,6,7,8,9,10,11,12,15\}$ & No & $\{13,14\}$ \\
10 & $\{0,1,2,3,4,5,6,7,8,9,10,11,13,14,15,20\}$ & No & $\{12,16,17,18,19\}$ \\
11 & $\{0,1,2,3,4,5,6,7,8,9,10,11,12,13,14,16,19\}$ & No & $\{15,17,18\}$ \\
12 & $\{0,1,2,3,4,5,6,7,8,9,10,11,12,13,14,15,16,17,18\}$ & Yes & -- \\
13 & $\{0,1,2,3,4,5,6,7,8,9,10,11,12,13,14,15,16,17,18,20\}$ & No & $\{19\}$ \\
14 & $\{0,1,2,3,4,5,6,7,8,9,10,11,12,13,14,15,16,17,18,19,20\}$ & Yes & -- \\
15 & $\{0,1,2,3,4,5,6,7,8,9,10,11,12,13,14,15,16,17,18,19,21,22\}$ & No & $\{20\}$ \\
16 & $\{0,1,2,3,4,5,6,7,8,9,10,11,12,13,14,15,16,17,18,19,20,22,23\}$ & No & $\{21\}$ \\
\bottomrule
\end{tabular}
\caption{Complete valuation spectra for best known sets of size $k=1,\dots,16$.}
\label{tab:spectra}
\end{table}

\subsection{Additional Constructions}

\begin{table}[ht]
\centering
\begin{tabular}{c|c|c|c}
\toprule
$k$ & Construction & $f(A)$ & Spectrum \\
\midrule
Any & Powers of two: $\{1,2,4,\dots,2^{k-1}\}$ & $k$ & $[0, k-1]$ \\
Any & Odd numbers: $\{1,3,5,\dots,2k-1\}$ & $\lfloor k/2 \rfloor + 2$ & $\{0,1,\dots,\lfloor k/2 \rfloor+1\}$ \\
3 & $\{1,3,5\}$ & 4 & $[0,3]$ \\
4 & $\{1,3,5,7\}$ & 5 & $[0,4]$ \\
5 & $\{1,2,5,9,15\}$ & 6 & $[0,5]$ \\
6 & $\{1,7,11,14,15,16\}$ & 7 & $[0,6]$ \\
8 & Greedy: $\{780,1097,1866,3083,7935,8620,9086,9773\}$ & 12 & $[0,11]$ \\
12 & Greedy: $\{6309,9914,12404,12666,15510,17728,17734,24752,25484,28688\}$ & 19 & $[0,18]$ \\
\bottomrule
\end{tabular}
\caption{Some explicit constructions and their valuation diversity.}
\label{tab:constructions}
\end{table}

\end{document}